\chapter{Pitfalls and Known Issues}\label{app:pitfalls}

Every item below was re-checked against this checkout by \file{examples/ex15_pitfalls.py};
its raw output is reproduced at the end of this appendix.

\begin{enumerate}
  \item \textbf{\code{from pystencils import *} fails} with an \code{AttributeError}:
        \code{__all__} lists \code{inspect_ast}, which does not exist.
  \item \textbf{\code{ps.show_code(kernel)} prints IR, not C} (it is deprecated and calls
        \code{ps.inspect} without \code{show_cpp=True}). Use \code{kernel.get_c_code()}.
  \item \textbf{The deprecated \code{CreateKernelConfig(cpu_openmp=True)} disables OpenMP.}
        Use \code{cfg.cpu.openmp.enable = True}.
  \item \textbf{\code{Field.create_fixed_size(..., memory_strides=...)} ignores the
        strides} and computes them from the layout instead.
  \item \textbf{Staggered kernels are not implemented in 2.0}:
        \code{ps.create_staggered_kernel} raises \code{NotImplementedError}.
  \item \textbf{\code{@ps.kernel} with \code{a if c else b} fails on Python 3.14}
        (\code{ast.NameConstant} was removed). Use \code{sp.Piecewise}.
  \item \textbf{Keyword arguments of \code{create_kernel} must be top-level options}:
        \code{create_kernel(asm, openmp=True)} raises a \code{TypeError}.
  \item \textbf{The default JIT ignores the kernel's target} and compiles with
        \code{-march=native} (\cref{sec:jit}).
  \item \textbf{\code{assume_inner_stride_one} is not checked at run time}: a strided
        NumPy view is accepted and gives \emph{wrong results} without an error.
  \item \textbf{CUDA warp-level reductions need an explicit \code{gpu.warp_size}}: with
        \code{assume_warp_aligned_block_size} alone, no \code{__shfl_xor_sync} is
        emitted, because the driver passes the unset option (not the default 32) to the
        platform.
  \item \textbf{The loop-independence check is order dependent}: the in-place stencil
        \texttt{f[0,0]}~$\leftarrow$~\texttt{f[1,0]} is accepted (\cref{sec:constraints}).
  \item \textbf{\code{Target.GenericCPU} prints as \code{Target._CPU}}.
  \item \textbf{\code{simplificationfactory.create_simplification_strategy()} returns
        \code{None}} (missing \code{return}).
  \item \textbf{Setting \code{block_size} on a dynamic GPU launch configuration has no
        effect} and raises no error. Use \code{fit_block_size}/\code{trim_block_size}, or
        \code{gpu.manual_launch_grid = True}.
  \item \textbf{Vectorized Philox needs a matching index width}: a \code{float32}
        generator on \code{X86_AVX} requires \code{index_dtype="int32"}.
  \item \textbf{Stale tutorial notebooks}: the stored outputs in
        \file{docs/source/tutorials/*.ipynb} come from older versions and use
        deprecated APIs (\code{pystencils.session}, \code{cpu_openmp},
        \code{get_parameters()}).
\end{enumerate}

\outsnippet[caption={Raw output of \texttt{examples/ex15\_pitfalls.py}.},basicstyle=\ttfamily\scriptsize]{ex15_pitfalls.results}
